%\chapter[gretl.Model class hierarchy]{\protect\classpymodel hierarchy}\label{sec:classpymodel_hierarchy}
%\begin{landscape}
%	\begin{figure}
%		\begin{tikzpicture}[
%			box/.style={rectangle, draw, rounded corners, minimum width=4cm, minimum height=1cm, align=center},
%			level distance=2.5cm, sibling distance=5cm,
%			edge/.style={-{Latex}}
%			]
%
%			% Base class
%			\node[box] (GretlModel) {GretlModel};
%
%			% NSE branch
%			\node[box, below=of GretlModel, xshift=-8cm] (GretlModel_NSE) {GretlModel\_NSE};
%			\node[box, below=of GretlModel_NSE, xshift=-2cm] (GretlModel_NSE_Driver) {GretlModel\_NSE\_Driver};
%			\node[box, below=of GretlModel_NSE_Driver, xshift=-3cm] (GretlModel_NSE_Driver_quantreg) {GretlModel\_NSE\_Driver\_quantreg};
%			\node[box, below=of GretlModel_NSE_Driver, xshift=3cm] (GretlModel_NSE_Driver_dpanel) {GretlModel\_NSE\_Driver\_dpanel};
%			\node[box, below=of GretlModel_NSE_Driver_dpanel] (GretlModel_NSE_Driver_tobit) {GretlModel\_NSE\_Driver\_tobit};
%
%			% NSE plugin
%			\node[box, below=of GretlModel_NSE, xshift=6cm] (GretlModel_NSE_nl) {GretlModel\_NSE\_nl};
%
%			% VAR branch
%			\node[box, below=of GretlModel, xshift=6cm] (GretlModel_VAR) {GretlModel\_VAR};
%			\node[box, below=of GretlModel_VAR] (GretlModel_VAR_VECM) {GretlModel\_VAR\_VECM};
%
%			% NME branch
%			\node[box, below=of GretlModel, yshift=-10cm] (GretlModel_NME) {GretlModel\_NME};
%
%			% Connections
%			\draw[edge] (GretlModel) -- (GretlModel_NSE);
%			\draw[edge] (GretlModel_NSE) -- (GretlModel_NSE_Driver);
%			\draw[edge] (GretlModel_NSE_Driver) -- (GretlModel_NSE_Driver_quantreg);
%			\draw[edge] (GretlModel_NSE_Driver) -- (GretlModel_NSE_Driver_dpanel);
%			\draw[edge] (GretlModel_NSE_Driver) -- (GretlModel_NSE_Driver_tobit);
%			\draw[edge] (GretlModel_NSE) -- (GretlModel_NSE_nl);
%			
%			\draw[edge] (GretlModel) -- (GretlModel_VAR);
%			\draw[edge] (GretlModel_VAR) -- (GretlModel_VAR_VECM);
%
%			\draw[edge] (GretlModel) -- (GretlModel_NME);
%
%		\end{tikzpicture}
%		\caption{\protect\classpymodel hierarchy}
%		\label{fig:model_class_hierarchy}
%	\end{figure}
%\end{landscape}


\chapter[gretl.Model class hierarchy]{\protect\classpymodel hierarchy}
\label{sec:classpymodel_hierarchy}

The \classpymodel class provides the common interface for
econometric models in \gretlpy. Model-specific classes inherit
from this base class, either directly or through intermediate
classes.

Figure~\ref{fig:model_class_hierarchy} illustrates the Python
class hierarchy. The three main branches are
\texttt{Model\_NSE}, \texttt{Model\_var}, and \texttt{System}.
The \texttt{Model\_NSE} branch contains single-equation models,
while \texttt{Model\_vecm} extends \texttt{Model\_var}.
The classes \texttt{Model\_ols\_add} and
\texttt{Model\_ols\_omit} inherit from \texttt{Model\_ols}.

\begin{landscape}
	\begin{figure}
		\centering
		
		\begin{tikzpicture}[
			font=\small,
			model/.style={
				rectangle,
				draw,
				rounded corners=2pt,
				minimum width=3.0cm,
				minimum height=0.62cm,
				align=center,
				inner sep=3pt
			},
			base/.style={
				model,
				fill=gray!20,
				font=\small\bfseries
			},
			special/.style={
				model,
				fill=blue!8
			},
			group/.style={
				draw=gray!65,
				rounded corners=5pt,
				fill=gray!4,
				inner sep=12pt
			},
			inheritance/.style={
				-{Latex[length=2mm]},
				thick
			}
			]
			
			% ==================================================
			% Main hierarchy
			% ==================================================
			
			\node[base] (base) at (0,0) {Model};
			
			\node[base] (nse) at (-7,-1.5) {Model\_NSE};
			\node[model] (var) at (0,-1.5) {Model\_var};
			\node[model] (system) at (7,-1.5) {System};
			
			\coordinate (rootsplit) at ($(base.south)+(0,-0.55)$);
			
			\draw[thick] (base.south) -- (rootsplit);
			\draw[inheritance] (rootsplit) -| (nse.north);
			\draw[inheritance] (rootsplit) -- (var.north);
			\draw[inheritance] (rootsplit) -| (system.north);
			
			% VAR branch
			
			\node[model] (vecm) at (0,-2.8) {Model\_vecm};
			
			\draw[inheritance] (var.south) -- (vecm.north);
			
			% ==================================================
			% Single-equation models
			% ==================================================
			
			\begin{scope}[shift={(-10,-5)}]
				
				% Column 1: OLS family
				
				\node[special] (ols) at (0,0) {Model\_ols};
				
				\node[special] (olsadd) at (-1.6,-1.15)
				{Model\_ols\_add};
				
				\node[special] (olsomit) at (1.6,-1.15)
				{Model\_ols\_omit};
				
				\node[model] (wls) at (0,-2.15) {Model\_wls};
				\node[model] (logit) at (0,-3.05) {Model\_logit};
				\node[model] (probit) at (0,-3.95) {Model\_probit};
				\node[model] (poisson) at (0,-4.85) {Model\_poisson};
				\node[model] (negbin) at (0,-5.75) {Model\_negbin};
				\node[model] (biprobit) at (0,-6.65) {Model\_biprobit};
				\node[model] (duration) at (0,-7.55) {Model\_duration};
				
				% Column 2
				
				\node[model] (heckit) at (5.7,0) {Model\_heckit};
				\node[model] (logistic) at (5.7,-0.9) {Model\_logistic};
				\node[model] (ar) at (5.7,-1.8) {Model\_ar};
				\node[model] (ar1) at (5.7,-2.7) {Model\_ar1};
				\node[model] (lad) at (5.7,-3.6) {Model\_lad};
				\node[model] (garch) at (5.7,-4.5) {Model\_garch};
				\node[model] (mpols) at (5.7,-5.4) {Model\_mpols};
				\node[model] (panel) at (5.7,-6.3) {Model\_panel};
				
				% Column 3
				
				\node[model] (ivreg) at (9.5,0) {Model\_ivreg};
				\node[model] (hsk) at (9.5,-0.9) {Model\_hsk};
				\node[model] (intreg) at (9.5,-1.8) {Model\_intreg};
				\node[model] (midasreg) at (9.5,-2.7) {Model\_midasreg};
				\node[model] (quantreg) at (9.5,-3.6) {Model\_quantreg};
				\node[model] (tobit) at (9.5,-4.5) {Model\_tobit};
				\node[model] (dpanel) at (9.5,-5.4) {Model\_dpanel};
				\node[model] (arima) at (9.5,-6.3) {Model\_arima};
				
				% Column 4: nonlinear models
				
				\node[model] (nls) at (13.3,0) {Model\_nls};
				\node[model] (mle) at (13.3,-0.9) {Model\_mle};
				\node[model] (gmm) at (13.3,-1.8) {Model\_gmm};
				
			\end{scope}
			
			% ==================================================
			% OLS inheritance
			% ==================================================
			
			\coordinate (olssplit) at ($(ols.south)+(0,-0.35)$);
			
			\draw[thick] (ols.south) -- (olssplit);
			\draw[inheritance] (olssplit) -| (olsadd.north);
			\draw[inheritance] (olssplit) -| (olsomit.north);
			
			% ==================================================
			% NSE group
			% ==================================================
			
			\begin{scope}[on background layer]
				
				\node[group,
				fit=(ols)(olsadd)(olsomit)
				(wls)(logit)(probit)(poisson)(negbin)
				(biprobit)(duration)
				(heckit)(logistic)(ar)(ar1)(lad)
				(garch)(mpols)(panel)
				(ivreg)(hsk)(intreg)(midasreg)
				(quantreg)(tobit)(dpanel)(arima)
				(nls)(mle)(gmm)
				] (nsegroup) {};
				
			\end{scope}
			
			% Group title: centered on upper border
			
			\node[
			font=\small\bfseries,
			fill=white,
			inner sep=3pt,
			anchor=south
			] at ($(nsegroup.north)+(0,0.08)$)
			{Subclasses of Model\_NSE};
			
			% Connection from Model_NSE to group
			
			\draw[inheritance]
			(nse.south) -- (nse.south |- nsegroup.north);
			
		\end{tikzpicture}
		
		\caption{Python class hierarchy of econometric models
			in \gretlpy. The shaded area groups subclasses of
			\texttt{Model\_NSE}. All classes within this group
			inherit directly from \texttt{Model\_NSE}, except
			\texttt{Model\_ols\_add} and \texttt{Model\_ols\_omit},
			which inherit from \texttt{Model\_ols}.}
		\label{fig:model_class_hierarchy}
		
	\end{figure}
\end{landscape}
